\section{Submodules}\label{sec:10-modules:04-submodules:1}

A \textbf{submodule} of \(M\) is a subset closed under \(+\), containing \(0\), and closed under the scalar action. In Lean, ``a subset closed under some operations'' is naturally a \texttt{structure} bundling a \emph{predicate} on \texttt{M} together with closure proofs. This is the same ``data + proofs'' idea as \texttt{Group}, just with the ``data'' being a property instead of an operation:

\begin{lstlisting}
structure Submodule {R : Type} (Rg : Ring R) {M : Type} (Mod : Module R Rg M) where
  carrier : M (*$\rightarrow$*) Prop
  zero_mem : carrier Mod.addGrp.toGroup.id
  add_mem : (*$\forall$*) {m n : M}, carrier m (*$\rightarrow$*) carrier n (*$\rightarrow$*) carrier (Mod.addGrp.op m n)
  smul_mem : (*$\forall$*) (r : R) {m : M}, carrier m (*$\rightarrow$*) carrier (Mod.smul r m)
\end{lstlisting}

\texttt{carrier : M → Prop} is the subset, viewed as its membership predicate (\texttt{carrier m} reads ``\texttt{m} is in the submodule''). This is the standard Lean idiom for ``subobject of \texttt{X}'': not a \texttt{Set X} wrapped separately, but a predicate directly, since \texttt{carrier m} \emph{is} a proposition usable directly as a hypothesis.

\begin{mathreading}

\texttt{Submodule Rg Mod} is a submodule \(N \le M\), presented by its membership predicate \(\chi_N : M \to \mathrm{Prop}\) (a subset \(N = \{\,m \in M \mid \chi_N(m)\,\}\)) together with closure proofs: \(0 \in N\), \(m,n \in N \Rightarrow m+n \in N\), and \(r\in R,\, m\in N
\Rightarrow r\cdot m \in N\). These are exactly the conditions for \(N\) to be an \(R\)-submodule: closed under the abelian-group operations and stable under scalars. So \(N\) inherits an \(R\)-module structure and the inclusion \(N \hookrightarrow M\) is \(R\)-linear. Categorically this is a \hyperref[sec:01-basics:04-terminology:1]{subobject} of \(M\) in \(R\text{-}\mathbf{Mod}\).

\end{mathreading}

\subsection{\texorpdfstring{Example: the submodule of even integers, as a \(\mathbb{Z}\)-submodule of \(\mathbb{Z}\)}{Example: the submodule of even integers, as a \textbackslash mathbb\{Z\}-submodule of \textbackslash mathbb\{Z\}}}\label{sec:10-modules:04-submodules:2}

The \(\mathbb{Z}\)-module structure on \(\mathbb{Z}\) itself, described only in passing so far, is required first:

\begin{lstlisting}
def intZModule : Module Int intRing Int where
  addGrp := intCommGroup
  smul := fun r m => r * m
  smul_add := by
    intro r m n
    exact Int.mul_add r m n
  add_smul := by
    intro r s m
    exact Int.add_mul r s m
  smul_smul := by
    intro r s m
    exact Int.mul_assoc r s m
  one_smul := by
    intro m
    exact Int.one_mul m
\end{lstlisting}

\begin{lstlisting}
-- NOTE: `intZModule.addGrp.op`/`.id`/`.smul` are definitionally (but not
-- syntactically) equal to plain `+`/`0`/`*` on `Int` (*\textemdash{}*) true by `rfl`, but
-- `rw` only fires on syntactic matches, so a bare `rw [...]` (or `ring`,
-- which this book doesn't import from Mathlib anyway) would leave
-- an unsolved goal even though the underlying statement is true. `show`
-- restates the goal in its reduced (defeq) form explicitly first.
def evenSubmodule : Submodule intRing intZModule where
  carrier := fun m => (*$\exists$*) k : Int, m = 2 * k
  zero_mem := (*$\langle$*)0, rfl(*$\rangle$*)
  add_mem := by
    intro m n (*$\langle$*)k, hk(*$\rangle$*) (*$\langle$*)j, hj(*$\rangle$*)
    refine (*$\langle$*)k + j, ?_(*$\rangle$*)
    show m + n = 2 * (k + j)
    rw [hk, hj, Int.mul_add]
  smul_mem := by
    intro r m (*$\langle$*)k, hk(*$\rangle$*)
    refine (*$\langle$*)r * k, ?_(*$\rangle$*)
    show r * m = 2 * (r * k)
    rw [hk, (*$\leftarrow$*) Int.mul_assoc, Int.mul_comm r 2, Int.mul_assoc]
\end{lstlisting}

Each closure proof follows the same shape: destructure the hypothesis to expose its witness (\texttt{⟨k, hk⟩} denotes ``\texttt{m} is even, with witness \texttt{k} and proof \texttt{hk : m = 2 * k}''), then supply a new witness and discharge the resulting \texttt{Int} equation. \texttt{zero\_\allowbreak{}mem}'s witness \texttt{0} makes the goal \texttt{id = 2 * 0}, true by \texttt{rfl} directly (both sides compute to \texttt{0}). \texttt{add\_\allowbreak{}mem} and \texttt{smul\_\allowbreak{}mem} require \texttt{show} first to reveal the goal's \texttt{+}/\texttt{*}-form before \texttt{rw} can find anything to rewrite, since \texttt{Mod.addGrp.op}/\texttt{Mod.smul} do not display as plain \texttt{+}/\texttt{*}, even though they compute to exactly that here.

\begin{mathreading}

This is the submodule \(2\mathbb{Z} \le \mathbb{Z}\) of even integers, \(2\mathbb{Z} = \{\, m \mid \exists k,\ m = 2k \,\}\). It's the image of multiplication-by-\(2\), and equivalently the principal ideal \((2)\trianglelefteq \mathbb{Z}\) (submodules of the \(\mathbb{Z}\)-module \(\mathbb{Z}\) are exactly its ideals). The closure checks are the elementary facts \(0 = 2\cdot 0\), \(2k + 2j = 2(k+j)\), and \(r\cdot(2k) = 2(rk)\), each an identity in the commutative ring \(\mathbb{Z}\). The existential witnesses \(0,\ k+j,\ rk\) are precisely the elements that show closure.

\end{mathreading}

\textbf{Mathlib equivalent.} Mathlib builds the submodule generated by a set \emph{generically}, using \href{https://loogle.lean-lang.org/?q=Submodule.span}{\texttt{Submodule.span}}. Thus \texttt{evenSubmodule}'s three hand-written closure proofs (\texttt{zero\_\allowbreak{}mem}/\texttt{add\_\allowbreak{}mem}/\texttt{smul\_\allowbreak{}mem}) become a single call, applicable to any generating set at all, not just \texttt{\{2\}}:

\begin{lstlisting}
def evenSubmodule' : Submodule Int Int := Submodule.span Int {(2 : Int)}

example : (2 : Int) (*$\in$*) evenSubmodule' := Submodule.subset_span rfl
\end{lstlisting}

\texttt{Submodule.span R s} is, by construction, the smallest submodule containing \texttt{s}. \href{https://loogle.lean-lang.org/?q=Submodule.subset_span}{\texttt{Submodule.subset\_\allowbreak{}span}} is the one fact that does all the work here: ``the generators are always members'', proved once, generically, instead of the book's custom \texttt{zero\_\allowbreak{}mem}/\texttt{add\_\allowbreak{}mem}/\texttt{smul\_\allowbreak{}mem} triple for this specific subset.
